The results demonstrate that a continuous Gaussian HMM trained via Baum-Welch achieves high distributional fidelity (90--94\% KS) and competitive temporal fidelity (ACF-MAE of 0.046--0.049) across all state resolutions tested, without requiring jump augmentation.
The following paragraphs interpret the key findings, address the Ryd\'{e}n et al.\ consensus, and identify limitations.

\paragraph{Resolving the Ryd\'{e}n et al.\ Limitation.}
The most significant finding of this study is that the CHMM achieves ACF-MAE of 0.046--0.049 on absolute returns, comparable to GARCH(1,1) (0.047), directly contradicting the widely cited conclusion of \citet{ryden1998stylized} that Gaussian HMMs cannot reproduce the slow decay of volatility autocorrelation.

The mechanism behind this contradiction is straightforward.
\citet{ryden1998stylized} tested only $K = 2$--$3$ states.
With two states, the sojourn time in each state follows a geometric distribution with a single rate parameter, producing a single exponential decay mode.
If that rate is fast (short mean sojourn), the ACF decays too quickly; if slow, the model spends too much time in one regime, distorting the marginal distribution.
This tradeoff is unresolvable at $K = 2$.

At $K \geq 3$ with quantile-based initialization, the Baum-Welch algorithm learns a transition matrix with heterogeneous diagonal entries.
Central (low-volatility) states develop high self-transition probabilities because calm market conditions are persistent, while tail states maintain lower self-transition probabilities reflecting the transient nature of extreme events.
The ACF of the resulting process is a weighted sum of $K$ exponential decay terms:
\begin{equation}
    \rho_{|G|}(\tau) \approx \sum_{k=1}^{K} w_k \, \lambda_k^{\tau}
    \label{eq:acf_mixture}
\end{equation}
where $\lambda_k$ are the eigenvalues of $\mathbf{T}$ (excluding the unit eigenvalue) and $w_k$ are weights determined by the emission parameters.
With $K \geq 6$ states spanning a range of persistence levels, this sum of exponentials can approximate the slow, quasi-hyperbolic ACF decay observed empirically.
This is, in effect, the same mathematical mechanism by which hidden semi-Markov models \citep{bulla2006stylized} achieve ACF fidelity---but realized within the standard HMM framework through increased state resolution rather than modified sojourn-time distributions.

It is notable that even $K = 3$ achieves ACF-MAE of 0.046 with our initialization, while \citet{ryden1998stylized} reported ACF failure at the same state count.
The difference is the initialization: the quantile-based strategy ensures that each state captures a distinct volatility regime from the start, guiding the EM algorithm toward a solution with differentiated emission parameters.
Random or uniform initialization---as was standard in the 1990s literature---increases the likelihood of convergence to degenerate local optima where states overlap, reducing the effective state resolution and reproducing the Ryd\'{e}n failure.

\paragraph{The Temporal-Distributional Tradeoff.}
The five-model comparison reveals that no single model dominates all seven metrics, confirming the temporal-distributional tradeoff identified in the discrete framework of \citet{alswaidan2026hybrid}.
GARCH(1,1) achieves the best ACF-MAE (0.047) but catastrophically fails KS (13\%), while the bootstrap achieves perfect KS (99.9\%) but the worst ACF-MAE (0.062) among non-trivial models.

The CHMM occupies a distinctive position in this tradeoff: it achieves KS (94.4\%) and AD (85.3\%) pass rates that approach the non-parametric bootstrap, while simultaneously matching GARCH's ACF-MAE (0.048 vs.\ 0.047).
It also achieves the best Wasserstein-1 (0.11) and Hellinger (0.077) distances among parametric models, and 100\% quantile coverage.
No other model in the comparison achieves this combination of distributional and temporal fidelity.

The GARCH distributional failure deserves specific attention.
Despite producing realistic conditional variance dynamics (high $\alpha + \beta$ near unity), the GARCH(1,1) model's Gaussian innovations generate a marginal distribution that deviates systematically from the observed heavy-tailed distribution.
The 13\% IS KS pass rate (and 3.2\% AD) indicates that the vast majority of GARCH-simulated paths are statistically distinguishable from the observed data by standard distributional tests.
This is not a calibration failure---the GARCH parameters are estimated by MLE---but a structural limitation of the single-regime conditional heteroscedasticity framework.

\paragraph{The Kurtosis Gap.}
Despite the CHMM's strong performance on distributional tests, a systematic kurtosis gap persists: the observed excess kurtosis is 7.71, while the CHMM produces 3.7--5.1 depending on $K$.
This gap is inherent to the Gaussian emission assumption.
A mixture of $K$ Gaussian distributions produces a marginal distribution whose excess kurtosis is bounded by the variance of the component means relative to the component variances.
For the kurtosis of the mixture to match the observed value of 7.71, the tail states would need substantially larger emission variances or more extreme emission means than the Baum-Welch algorithm selects---but such parameters would distort the central body of the distribution and reduce KS pass rates.

The kurtosis gap explains why the AD pass rate (85.3\%) is consistently lower than the KS pass rate (94.4\%): the Anderson--Darling test assigns greater weight to the tails, where the Gaussian mixture's lighter-than-observed tails produce the largest discrepancies.
Replacing Gaussian emissions with Student-$t$ distributions---whose heavier tails can produce arbitrary excess kurtosis---would address this gap while retaining the Baum-Welch estimation framework.
The M-step update for Student-$t$ emissions requires iterative estimation of the degrees-of-freedom parameter $\nu_k$ per state, adding computational cost but no conceptual difficulty.

\paragraph{Why Jumps Are Unnecessary.}
The discrete framework of \citet{alswaidan2026hybrid} required a Poisson jump-duration mechanism to force the model into tail states for empirically realistic durations.
The continuous model achieves comparable performance without this mechanism.
The explanation lies in the richer parameterization of continuous emissions.

In the discrete model, the emission distributions were estimated \emph{separately} from the transition matrix: observations were first binned into quantile-defined states, then the transition matrix was estimated by frequency counting, and finally within-bin Student-$t$ distributions were fitted.
This three-stage procedure meant that the transition matrix was not optimized jointly with the emissions---so the learned $\mathbf{T}$ might not produce sufficient tail-state persistence, necessitating the jump mechanism as a corrective.

In the continuous model, the Baum-Welch algorithm jointly optimizes $\mathbf{T}$ and $\{(\mu_k, \sigma_k)\}$.
The emission parameters and transition probabilities co-adapt: if a tail state has a distinctive emission distribution that captures extreme observations well, the posterior $\gamma_t(k)$ will assign high probability to that state during extreme events, and the M-step will increase the self-transition probability $T_{kk}$ to reflect the clustering of such events.
This joint optimization is the key advantage of the continuous approach---it allows the model to learn the persistence of tail regimes directly from data, rather than imposing it externally.

The elimination of the jump mechanism removes two hyperparameters ($\epsilon$, $\lambda$) and the grid search required to calibrate them.
For the discrete model, this grid search evaluated $6 \times 8 = 48$ parameter combinations, each requiring 200 simulated paths, for a total of 9,600 simulations per $K$ value.
The continuous model requires only the Baum-Welch iterations (20--40 to convergence), making it substantially more computationally efficient.

\paragraph{State Resolution as the Key Variable.}
The state resolution analysis (Table~\ref{tab:sensitivity}) shows that $K$ is the single most important design choice for the CHMM.
KS pass rates improve from 90.1\% ($K = 3$) to 94.2\% ($K = 11$), while ACF-MAE remains stable and quantile coverage is uniformly 100\%.
The practical implication is that practitioners should choose $K$ based on the distributional fidelity required: $K = 3$ suffices for applications that need approximate regime identification with reasonable ACF behavior, while $K = 11$--$13$ maximizes distributional accuracy.

The stability of ACF-MAE across $K$ is theoretically predicted by the eigenvalue structure of $\mathbf{T}$.
The dominant eigenvalue $\lambda_2$ (which controls the slowest ACF decay mode) depends on the maximum self-transition probability in the matrix.
At all $K$ values, the Baum-Welch algorithm learns at least one high-persistence central state (e.g., $T_{kk} > 0.9$), ensuring that $\lambda_2$ remains close to unity regardless of the total number of states.
Additional states refine the distributional shape of the mixture but do not substantially affect the dominant eigenvalue.

\paragraph{Cross-Asset Robustness and OoS Degradation.}
The cross-asset results confirm that the CHMM generalizes to equities with diverse risk profiles (NVDA: 97.4\% IS KS; JNJ: 99.8\%; JPM: 98.6\%).
The OoS degradation for JNJ (68.7\%) and JPM (79.8\%) reflects the stationarity assumption: the CHMM's transition matrix and emission parameters are fixed at their IS-estimated values, so regime shifts between the training and test periods---which may arise from company-specific events, sector rotations, or macroeconomic shocks---are not captured.
NVDA's remarkably high OoS KS (98.5\%) suggests that its 2025 dynamics remained within the distributional envelope learned from 2014--2024, possibly because the technology sector's volatility profile was more stable over this period.

A time-varying extension---either through rolling-window re-estimation or Bayesian online learning of the transition matrix---would address the stationarity limitation and likely improve OoS performance for assets with non-stationary regime dynamics.

\paragraph{Limitations.}
Several limitations qualify the conclusions of this study:
\begin{enumerate}
    \item \textit{Gaussian emission kurtosis gap.}
    The Gaussian assumption systematically underestimates excess kurtosis (4.2 simulated vs.\ 7.7 observed at $K = 13$), limiting the model's ability to reproduce extreme tail events.
    Student-$t$ or generalized hyperbolic emissions would address this.

    \item \textit{Stationarity assumption.}
    The transition matrix and emission parameters are estimated from the full IS window and held constant.
    Structural breaks (e.g., the March 2020 COVID-19 crash) may not be adequately captured if the IS window does not contain comparable events.

    \item \textit{Single OoS window.}
    The out-of-sample evaluation covers a single 219-day window (2025).
    A rolling or expanding-window evaluation would provide more robust generalization assessment.

    \item \textit{KS test i.i.d.\ assumption.}
    The two-sample KS test assumes i.i.d.\ observations, but each simulated path exhibits temporal dependence by construction.
    Volatility clustering in the simulated paths may inflate pass rates slightly relative to a block-bootstrap calibration \citep{alswaidan2026hybrid}.

    \item \textit{Gaussian emission symmetry.}
    The Gaussian emission for each state is symmetric, while empirical returns exhibit negative skewness ($-0.75$ for SPY).
    The mixture can produce overall skewness through asymmetric state occupancies and means, but may not fully capture the within-regime asymmetry.

    \item \textit{No formal model selection.}
    We do not apply information criteria (AIC, BIC) for choosing $K$, instead reporting results across all tested values.
    A principled model selection procedure would provide guidance on the optimal state count.
\end{enumerate}
